\documentclass[11pt]{article}
\usepackage[a4paper,margin=1in]{geometry}
\usepackage{fontspec}
\usepackage[boldfont=false]{xeCJK}
\setCJKmainfont{FandolSong-Regular.otf}
\usepackage{amsmath}
\usepackage{amssymb}
\usepackage{array}
\usepackage{booktabs}
\usepackage{graphicx}
\usepackage{adjustbox}
\usepackage{enumitem}
\setlist[itemize]{topsep=2pt,partopsep=0pt,parsep=0pt,itemsep=2pt,leftmargin=1.4em}
\setlist[enumerate]{topsep=2pt,partopsep=0pt,parsep=0pt,itemsep=2pt,leftmargin=1.6em}
\usepackage{parskip}
\usepackage{fvextra}
\fvset{breaklines=true,breakanywhere=true,fontsize=\small}
\setlength{\parindent}{0pt}

\begin{document}

\begin{center}
  {\LARGE\bfseries Homeostasis}\\[0.35em]
  {\large A Level Biology}
\end{center}
\vspace{0.6em}

\section*{What homeostasis is}

\textbf{Homeostasis} 稳态 means keeping the conditions inside the body steady, even when the outside changes. Keeping things like temperature, water and blood glucose steady lets enzymes and cells work properly all the time.

\subsection*{The principles of homeostasis}

Most homeostasis follows the same plan:

\begin{itemize}
\item a change (an internal or external \textbf{stimulus} 刺激) is detected by a \textbf{receptor} 受体.

\item a coordination system carries the message — either the \textbf{nervous system} 神经系统 (using nerve signals) or the \textbf{endocrine system} 内分泌系统 (using \textbf{hormones} 激素).

\item an \textbf{effector} 效应器 (a muscle or a \textbf{gland} 腺体) makes a response that corrects the change.

\end{itemize}

This works by \textbf{negative feedback} 负反馈: a change away from the normal level triggers a response that pushes it back towards normal.

\section*{The liver and urea}

The body cannot store extra amino acids. In the \textbf{liver} 肝脏, the process of \textbf{deamination} 脱氨基作用 removes the amino group from excess \textbf{amino acids} 氨基酸, and this is turned into \textbf{urea} 尿素. The urea is carried in the blood to the kidneys to be removed.

\section*{The kidney}

The \textbf{kidney} 肾脏 cleans the blood and controls its water content. Its parts are:

\begin{itemize}
\item an outer \textbf{fibrous capsule},

\item an outer region, the \textbf{cortex} 皮质,

\item an inner region, the \textbf{medulla} 髓质,

\item a central space, the \textbf{renal pelvis} 肾盂, which collects urine,

\item the \textbf{ureter} 输尿管, which carries urine to the bladder,

\item branches of the renal artery (bringing blood in) and renal vein (taking blood out).

\end{itemize}

\subsection*{The nephron}

Each kidney holds about a million tiny tubes called \textbf{nephrons} 肾单位. Along a nephron are: the \textbf{glomerulus} 肾小球 (a knot of capillaries), the \textbf{Bowman's capsule} 鲍曼囊 around it, the \textbf{proximal convoluted tubule} 近曲小管, the \textbf{loop of Henle} 亨利环, the \textbf{distal convoluted tubule} 远曲小管, and the \textbf{collecting duct} 集合管.

\subsection*{Making urine}

\begin{enumerate}
\item \textbf{Ultrafiltration} 超滤 happens in the Bowman's capsule. The blood in the glomerulus is under high pressure, so water and small molecules (\textbf{glucose} 葡萄糖, \textbf{ions} 离子, urea) are pushed out into the capsule, forming the \textbf{filtrate} 滤液. Blood cells and large proteins are too big to pass, so they stay in the blood.

\item \textbf{Selective reabsorption} 选择性重吸收 happens in the proximal convoluted tubule. Useful substances are taken back into the blood. All the glucose and much of the water and ions are reabsorbed here. The tubule wall is well suited to this: its cells have \textbf{microvilli} 微绒毛 to give a large surface area, and many \textbf{mitochondria} 线粒体 to power \textbf{active transport} 主动运输.

\end{enumerate}

\section*{Osmoregulation}

\textbf{Osmoregulation} 渗透调节 controls the water content of the blood. It is run by the brain:

\begin{itemize}
\item the \textbf{hypothalamus} 下丘脑 detects the \textbf{water potential} 水势 of the blood.

\item when the blood is too concentrated, the \textbf{pituitary gland} 垂体 releases \textbf{antidiuretic hormone} 抗利尿激素 (ADH).

\item ADH makes the collecting ducts more permeable to water by adding water channels called \textbf{aquaporins} 水通道蛋白.

\item more water is then reabsorbed back into the blood, so less, more concentrated urine is made. This is negative feedback.

\end{itemize}

\section*{Controlling blood glucose by cell signalling}

When blood glucose falls, the hormone \textbf{glucagon} 胰高血糖素 is released. It shows how a hormone passes its message into a cell — \textbf{cell signalling} 细胞信号传递:

\begin{enumerate}
\item glucagon binds to a receptor on the liver cell surface, causing a \textbf{conformational change} 构象变化 (a change in the receptor's shape).

\item this activates a \textbf{G-protein} G蛋白, which switches on the enzyme \textbf{adenylyl cyclase} 腺苷酸环化酶.

\item adenylyl cyclase makes a \textbf{second messenger} 第二信使 inside the cell, called \textbf{cyclic AMP} 环腺苷酸 (cAMP).

\item cAMP activates \textbf{protein kinase A} 蛋白激酶A, which starts an \textbf{enzyme cascade} 酶级联反应 — one \textbf{enzyme} 酶 switches on the next, by \textbf{phosphorylation} 磷酸化.

\item because each enzyme switches on many of the next, the signal is greatly \textbf{amplified} 放大.

\item the final enzyme breaks down \textbf{glycogen} 糖原 into glucose, which raises the blood glucose level.

\end{enumerate}

\subsection*{Negative feedback and blood glucose}

Blood glucose is held steady by two hormones working against each other:

\begin{itemize}
\item when glucose is \textbf{high}, \textbf{insulin} 胰岛素 makes muscle and liver cells take in glucose and store it as glycogen, lowering the level.

\item when glucose is \textbf{low}, glucagon makes liver cells break glycogen back into glucose, raising the level.

\end{itemize}

\subsection*{Measuring glucose}

\textbf{Test strips} 试纸 and \textbf{biosensors} 生物传感器 measure glucose in blood or urine. They use the enzymes \textbf{glucose oxidase} 葡萄糖氧化酶 and \textbf{peroxidase} 过氧化物酶, which react with glucose to give a colour change (or an electric signal in a biosensor) that shows how much glucose is present.

\section*{Homeostasis in plants: the stomata}

\textbf{Stomata} 气孔 are pores in a leaf. The plant opens and closes them to balance two needs: letting in \textbf{carbon dioxide} 二氧化碳 for \textbf{photosynthesis} 光合作用, and limiting water loss by \textbf{transpiration} 蒸腾作用. Stomata usually open by day and close at night, following a daily rhythm.

Each stoma is opened and closed by two \textbf{guard cells} 保卫细胞 around it:

\begin{itemize}
\item to \textbf{open}: the guard cells pump in ions, so their water potential falls and water enters by \textbf{osmosis} 渗透. They swell and bend apart, opening the pore.

\item to \textbf{close}: ions leave, water follows out, the guard cells go floppy, and the pore closes.

\end{itemize}

When the plant is short of water (\textbf{water stress} 水分胁迫), the hormone \textbf{abscisic acid} 脱落酸 is released. It makes the guard cells lose ions and water so the stomata close, saving water. In this signalling, \textbf{calcium ions} 钙离子 act as a second messenger inside the guard cells.


\end{document}
